\documentclass[12pt]{article}
\usepackage[UTF8]{ctex}
\usepackage[margin=2.2cm]{geometry}
\usepackage{enumitem,booktabs,longtable,array,xcolor,hyperref}
\hypersetup{colorlinks=true,linkcolor=blue,urlcolor=blue}
\setlist{nosep,leftmargin=2em}
\title{日序 Daymark\個人行程與任務管理系統\\產品需求文件（PRD）}
\author{產品規劃草案}
\date{版本 1.0 \quad｜\quad 2026 年 11 月 6 日}
\begin{document}
\maketitle
\begin{center}\textit{以清晰、可持續的日常規劃，協助個人兼顧行程、待辦、習慣與長期目標。}\end{center}
\tableofcontents\newpage
\section{文件資訊}
\begin{tabular}{p{3cm}p{10cm}}\toprule
產品名稱 & 日序 Daymark（暫定）\\\midrule
文件版本 & 1.0（初步需求草案）\\
狀態 & 待利害關係人確認；功能細節尚有待訪談\\
目標使用者 & 個人使用者\\
平台方向 & 響應式網頁，支援手機與桌面瀏覽器\\
範圍聲明 & 本文件描述產品需求；目前 HTML/CSS/JavaScript 原型屬前端展示，不代表已具備雲端同步、郵件發送或排程推播能力。\\\bottomrule\end{tabular}
\section{產品概述}
\subsection{產品願景}
提供一個個人化、低負擔且易辨識的規劃空間，將有時間的行程、彈性待辦、獨立習慣追蹤，以及跨月／年度目標放在一致的體驗中。介面採明亮易讀的科幻霓虹視覺語言，避免深色背景造成辨識困難，並以較大的圖示與觸控目標優先考量手機操作。
\subsection{問題陳述}
使用者需要在不同時間尺度管理承諾：今天要做什麼、何時要做、習慣是否維持，以及幾個月或數年內的目標如何推進。若分散在不同工具，容易遺漏截止日、難以建立日常節奏，也難以看見目標進度。
\subsection{產品目標}
\begin{enumerate}
\item 在單一入口同時呈現有時段行程與未排時段待辦。
\item 讓任務期限、優先級及完成狀態容易建立與追蹤。
\item 以獨立追蹤方式記錄習慣，不要求與目標綁定。
\item 支援具期限與里程碑的中長期目標。
\item 規劃手機與桌面瀏覽器提醒，並以電子郵件作為另一提醒管道。
\end{enumerate}
\subsection{非目標（首版暫定）}
團隊協作、公開社群、專案管理依賴圖、AI 自動排程、支付／訂閱、原生手機 App、與外部行事曆雙向同步。若後續訪談確認有需求，再調整範圍。
\section{使用者與主要情境}
\textbf{主要使用者：}希望以個人方式安排每天、記錄待辦與截止時間、持續建立習慣，並追蹤月／年目標的人。\
\textbf{情境 A：每日規劃} 使用者開啟今日總覽，查看時間軸和未排定時間的待辦，完成後勾選任務。\
\textbf{情境 B：期限管理} 使用者建立有日期、可選時間、分類與優先級的任務，在截止日前收到提醒。\
\textbf{情境 C：習慣追蹤} 使用者每日對習慣打卡，查看本週紀錄與連續紀錄；習慣可獨立於目標管理。\
\textbf{情境 D：長期規劃} 使用者建立數月或年度目標，設定期限、里程碑，並透過完成里程碑更新進度。\
\textbf{情境 E：跨裝置提醒} 使用者在手機或桌面瀏覽器授權通知，並可選擇電子郵件提醒；多裝置偏好與資料同步需由後端支援。
\section{資訊架構與導覽}
\begin{longtable}{p{3cm}p{10cm}}\toprule
區域 & 內容與責任\\\midrule
今日總覽 & 今日日期、時間軸行程、今日待辦、每日習慣摘要與完成概況\\
行事曆 & 月檢視、日期導覽、該日行程／任務摘要\\
所有任務 & 任務清單、完成狀態與全部／未完成／已完成篩選\\
習慣追蹤 & 習慣列表、打卡、連續天數與週紀錄\\
長期目標 & 目標列表、期限、進度及可勾選里程碑\\
提醒與設定 & 瀏覽器通知、郵件提醒、習慣提醒與提醒提前時間\\\bottomrule\end{longtable}
手機版需有明確的工作區導覽入口；開啟側邊導覽後可切換上述區域，選擇區域後自動收合，遮罩或返回操作可關閉。主要操作目標需適合觸控並有清楚的焦點／選取狀態。
\section{功能需求}
\begin{longtable}{p{1.2cm}p{4cm}p{7.8cm}}\toprule
ID & 功能 & 需求描述與初步驗收條件\\\midrule
FR-01 & 響應式工作區 & 桌面呈現側邊導覽；手機提供可開啟／關閉的導覽抽屜，能進入所有工作區，且不被頁面內容遮擋。\\
FR-02 & 今日總覽 & 顯示所選日期的時間軸行程、未排時間待辦與習慣摘要；完成任務後即時更新數量／進度。\\
FR-03 & 建立／管理任務 & 任務至少包含名稱、日期；時間為選填；可選分類與優先級。可切換完成狀態並按狀態篩選。\\
FR-04 & 行事曆 & 月曆可前後切換月份、返回本月，並以日期格摘要呈現行程／任務。\\
FR-05 & 習慣追蹤 & 可建立習慣、記錄完成／未完成，呈現週紀錄與連續天數。習慣資料模型不得強制關聯目標。\\
FR-06 & 目標及里程碑 & 目標具有標題、期限與里程碑；勾選里程碑後可更新完成進度，並提示逾期／即將到期狀態（後端版待細定）。\\
FR-07 & 行程提醒 & 可設定行程開始前提醒，預設示例為 10 分鐘；使用者能調整或關閉。\\
FR-08 & 任務截止提醒 & 可在截止前通知，預設示例為 1 天前；需支援使用者調整提醒時間。\\
FR-09 & 習慣未完成提醒 & 使用者可設定晚間提醒時間；僅對當日尚未完成且符合頻率的習慣提醒。\\
FR-10 & 目標期限提醒 & 里程碑／目標到期前提醒，預設示例為 3 天前；支援開關及調整。\\
FR-11 & 提醒管道 & 支援瀏覽器通知及電子郵件。瀏覽器通知需取得授權；電子郵件地址驗證、退訂與寄送失敗處理列入後端設計。\\
FR-12 & 持久化與跨裝置同步 & 使用者操作資料需在重新登入／重新整理後保留；同一帳號在已支援裝置間同步。此為產品目標，現有原型尚未實作。\\
FR-13 & 空／錯誤狀態 & 無資料時提供清楚空狀態與新增入口；網路失敗、通知未授權或儲存失敗時提供可理解訊息及重試方式。\\\bottomrule\end{longtable}
\section{資料需求（初步）}
\begin{longtable}{p{2.3cm}p{10.7cm}}\toprule
實體 & 建議欄位\\\midrule
Task & id、使用者 id、標題、日期、開始時間（選填）、截止時間（選填）、分類、優先級、完成狀態、提醒偏好、建立／更新時間\\
Event & id、使用者 id、標題、開始／結束時間、地點（選填）、分類、重複規則（選填）、提醒偏好\\
Habit & id、使用者 id、名稱、頻率規則、目標量（選填）、啟用狀態\\
HabitLog & id、習慣 id、日期、完成狀態、紀錄時間、備註（選填）\\
Goal & id、使用者 id、標題、描述（選填）、開始日、期限、狀態、進度計算方式\\
Milestone & id、目標 id、標題、期限（選填）、完成狀態、完成時間、排序\\
ReminderPreference & 使用者 id、管道、類型、提前量／時刻、啟用狀態、時區\\\bottomrule\end{longtable}
\textbf{資料規則：}日期時間應保留使用者時區語意；習慣完成紀錄按日期管理；刪除目標時需明確定義里程碑的連動處理；目標進度採里程碑比例或手動設定，待產品決策確認。
\section{非功能需求}
\begin{enumerate}
\item \textbf{易讀性與可用性：}高對比文字、清晰狀態色、鍵盤可操作、焦點可見；顏色不可作為唯一狀態提示。霓虹色用於強調，不大面積降低內文對比。
\item \textbf{響應式：}支援常見手機直向與桌面視窗；手機操作適用觸控，避免內容水平溢出。
\item \textbf{效能：}核心頁面在一般網路環境快速可互動；避免不必要的大型前端依賴。
\item \textbf{可靠性：}儲存與通知結果必須可辨認；跨裝置同步失敗不得靜默遺失資料。
\item \textbf{安全與隱私：}個人行程、習慣、目標皆屬私人資料；需登入授權、傳輸加密、最小權限及帳戶資料匯出／刪除方案（上線前確認）。
\item \textbf{可及性：}按鈕具語意標籤、表單欄位有標籤、動態回饋可由輔助科技辨識，並尊重減少動態效果偏好。
\end{enumerate}
\section{提醒與同步流程（預期）}
任務／行程／習慣／里程碑設定提醒 → 後端依使用者時區計算觸發時間 → 透過已授權的瀏覽器推播或已驗證電子郵件發送 → 記錄發送／失敗結果 → 使用者可管理偏好與取消提醒。瀏覽器頁面常駐執行的前端計時器不足以保證跨裝置、關閉瀏覽器後的通知，因此正式服務需後端排程及推播／郵件服務；推播技術與郵件供應商尚未選定。
\section{衡量指標}
\begin{itemize}
\item 啟用：新使用者完成首次任務或建立第一個習慣的比例。
\item 日常價值：每週活躍天數、任務建立至完成比例、習慣打卡週頻率。
\item 長期價值：目標建立率、里程碑完成率、目標期限前進度。
\item 提醒品質：通知送達率、提醒後任務完成率、提醒關閉／退訂率。
\item 體驗品質：手機導覽成功率、表單完成率、錯誤率及可用性測試任務成功率。
\end{itemize}
指標目標值與分析事件需在試用基準資料取得後設定，避免未驗證的數字成為承諾。
\section{首版範圍與優先級}
\begin{description}
\item[P0（首版必要）] 響應式導覽、今日總覽、任務建立／完成／篩選、月曆瀏覽、習慣建立／打卡、目標／里程碑、基本資料持久化、重要空狀態。
\item[P1（首版或緊接版本）] 瀏覽器通知、郵件提醒、提醒偏好、多裝置同步與帳號登入、資料匯出。
\item[P2（未來評估）] 外部行事曆整合、重複行程、智慧排程、進階分析與客製化主題。
\end{description}
\section{目前原型與落差}
目前原型已有五個主要工作區、任務新增與完成操作、月曆月份切換、習慣打卡、目標里程碑互動、提醒設定介面，以及手機工作區選單和霓虹科幻樣式。原型資料存在頁面執行期間的 JavaScript 記憶體中；重新載入後重設為示範內容。行程／任務／習慣／目標資料尚未串接持久化服務，電子郵件與跨裝置同步僅為展示，瀏覽器通知目前僅有權限互動示意而非可靠排程推播。上述項目不可視為已交付的正式功能。
\section{待確認事項}
\begin{enumerate}
\item 首版是否需登入及雲端同步，或先以單一裝置本機儲存驗證？
\item 行程是否需支援重複、結束時間、地點及外部行事曆匯入／同步？
\item 目標進度由里程碑比例自動計算，還是允許使用者手動更新？
\item 習慣頻率要支援每日、每週指定日及自訂次數嗎？補打卡規則為何？
\item 郵件提醒是否需要每日摘要？可用哪個電子郵件地址或帳戶管理方式？
\item 提醒是否可對每一筆項目個別設定？勿擾時段與時區變更如何處理？
\item 個人資料保留、匯出及帳號刪除的政策與需求為何？
\item 首版裝置／瀏覽器支援矩陣與可及性標準需採何種基準？
\end{enumerate}
\section{假設、限制與風險}
\begin{itemize}
\item 本文件根據目前對話與前端原型整理，尚未完成完整需求訪談或技術評估。
\item 跨裝置同步、電子郵件及背景推播需要後端、認證、排程與第三方服務；成本、平台限制與隱私條款需另行評估。
\item 習慣連續紀錄可能使使用者感到壓力，設計上應採鼓勵式語氣，容許中斷後重新開始。
\item 霓虹視覺若對比不足會影響可讀性，應以可讀性與可及性為優先，並透過裝置實測調校。
\end{itemize}
\section{驗收原則}
首版驗收至少應包含：核心流程可在手機與桌面完成；新資料能持久保存並於登入裝置間同步（若納入首版）；提醒只在使用者啟用與授權後送出且時區正確；資料狀態更新一致；空狀態、失敗狀態、鍵盤導覽及文字對比通過產品／QA 檢核。
\end{document}